% Diagramme de cas d'utilisation — CIVCO BTP (TikZ)
\begin{figure}[H]
\centering
\resizebox{\textwidth}{!}{%
\begin{tikzpicture}[
  actor/.style={font=\small\bfseries, align=center},
  uc/.style={ellipse, draw=ensamBlue, fill=ensamBlue!8, minimum width=2.6cm,
             minimum height=0.85cm, align=center, font=\scriptsize},
  sys/.style={rectangle, draw=ensamBlue!80, dashed, rounded corners=6pt},
  line/.style={draw=gray!60}
]
  % Cas d'utilisation — définis en premier
  \node[uc] (uc1)  at (-2.8, 6.0)  {S'authentifier};
  \node[uc] (uc2)  at (0.0, 6.0)   {Consulter le\\tableau de bord};
  \node[uc] (uc3)  at (2.8, 6.0)   {Gérer les\\utilisateurs};
  \node[uc] (uc4)  at (-3.2, 4.2)  {Gérer les\\projets};
  \node[uc] (uc5)  at (0.0, 4.2)   {Gérer les\\clients};
  \node[uc] (uc6)  at (3.2, 4.2)   {Gérer les\\tâches};
  \node[uc] (uc7)  at (-3.2, 2.4)  {Créer / suivre\\les devis};
  \node[uc] (uc8)  at (0.0, 2.4)   {Émettre les\\factures};
  \node[uc] (uc9)  at (3.2, 2.4)   {Produire les\\bons de livraison};
  \node[uc] (uc10) at (-3.2, 0.6)  {Gérer les\\contrats};
  \node[uc] (uc11) at (0.0, 0.6)   {Consulter la\\carte des projets};
  \node[uc] (uc12) at (3.2, 0.6)   {Configurer\\l'application};
  \node[uc] (uc13) at (-1.6,-1.2)  {Consulter\\l'historique};
  \node[uc] (uc14) at (1.6,-1.2)   {Échanger via\\la messagerie};
  \node[uc] (uc15) at (4.8,-0.8)   {Consulter projets\\et documents};

  % Frontière système (après les nœuds)
  \node[sys, fit=(uc1)(uc2)(uc3)(uc4)(uc5)(uc6)(uc7)(uc8)(uc9)(uc10)(uc11)(uc12)(uc13)(uc14)(uc15),
        inner sep=0.5cm,
        label={[font=\small\bfseries]above:Application de gestion CIVCO BTP}] (system) {};

  % Acteurs
  \node[actor] (admin)  at (-6.8, 5.2)  {Administrateur};
  \node[actor] (pm)     at (-6.8, 3.4)  {Chef de\\projet};
  \node[actor] (com)    at (-6.8, 1.6)  {Commercial};
  \node[actor] (acc)    at (-6.8,-0.2)  {Comptable};
  \node[actor] (user)   at (-6.8,-2.0)  {Utilisateur\\interne};
  \node[actor] (client) at (6.8, 1.0)   {Client\\externe};

  % Liaisons acteurs
  \foreach \a in {admin,pm,com,acc,user} {
    \draw[line] (\a.east) -- (uc1.west);
    \draw[line] (\a.east) -- (uc2.west);
  }
  \draw[line] (admin.east) -- (uc3.west);
  \draw[line] (pm.east) -- (uc4.west);
  \draw[line] (pm.east) -- (uc6.west);
  \draw[line] (com.east) -- (uc5.west);
  \draw[line] (com.east) -- (uc7.west);
  \draw[line] (acc.east) -- (uc8.west);
  \draw[line] (acc.east) -- (uc9.west);
  \draw[line] (pm.east) -- (uc10.west);
  \draw[line] (user.east) -- (uc6.west);
  \draw[line] (admin.east) -- (uc11.west);
  \draw[line] (admin.east) -- (uc12.west);
  \draw[line] (admin.east) -- (uc13.west);
  \draw[line] (pm.east) -- (uc14.west);
  \draw[line] (client.west) -- (uc15.east);
  \draw[line] (client.west) -- (uc7.east);
\end{tikzpicture}%
}
\caption{Diagramme de cas d'utilisation de l'application CIVCO BTP}
\label{fig:uc_diagram}
\end{figure}
